\subsection{Chebyshev parametrizations and the scope of the result}

Let $T_n$ denote the classical Chebyshev polynomial,
$T_n(\cos\theta)=\cos(n\theta)$, and use the normalization
$C_n(t)=2T_n(t/2)$. Thus
\[
\begin{aligned}
C_3&=t^3-3t, & C_4&=t^4-4t^2+2,\\
C_5&=t^5-5t^3+5t, & C_7&=t^7-7t^5+14t^3-7t.
\end{aligned}
\]
The real harmonic knots $H(3,4,5)$ and $H(3,5,7)$ are the trefoil and
the figure-eight, respectively; see \cite[Introduction and
Definition~1.2]{KP}. The change from $T_n$ to $C_n$ is a linear
reparametrization together with a common linear scaling of the ambient
coordinates.

The normalized Chebyshev trefoil is already covered by the construction:
\[
(C_3(t),C_4(t),C_5(t))=M(\gamma_T(t)),\qquad
M(x,y,z)=(x,y+2,z-5x).
\]
Consequently its rectifier and inverse are
\[
\Psi=\Phi_T\comp M^{-1},\qquad
\Psi^{-1}=M\comp\Phi_T^{-1},\qquad
M^{-1}(x,y,z)=(x,y-2,z+5x).
\]
In particular, $\Psi(C_3(t),C_4(t),C_5(t))=(t,0,0)$ over
$K=\Q(\sqrt{-3})$. This is a direct consequence of the explicit affine
normalization, rather than an additional hypothesis.

The harmonic model treated above is exactly $(C_3,C_5,C_7)$. The
polynomial identity
\[
C_3C_7-C_5^2+2=C_4
\]
explains the modification
$\beta(x,y,z)=(x,xz-y^2+2,y)$ connecting it to $(C_3,C_4,C_5)$.
The map $\beta$ becomes invertible after inverting $x$; the transfer
construction and its two exact divisions produce the global polynomial
automorphism.

These two cases do not establish rectifiability for all Chebyshev or
harmonic parametrizations. A further application must supply a closed
complex embedding, a coordinate restricting to its parameter, and the
hypotheses of the transfer theorem. Determining which additional
Chebyshev parametrizations satisfy these conditions remains a natural
question. In particular, the real knot classification in \cite{KP}
does not by itself establish polynomial equivalence of their complex
embeddings.
